\documentclass{article}

\usepackage[utf8]{inputenc}
\usepackage[T1]{fontenc}
\usepackage{amsmath,amssymb,amsthm}
\usepackage{mathtools}
\usepackage{hyperref}
\usepackage{booktabs}

% Lightweight citation command for this standalone draft (no natbib).
% Renders \citep{key} as (Paper N) based on the key.
\makeatletter
\newcommand{\citep}[2][]{%
  \begingroup
  \def\@tmpkey{#2}%
  \ifx\@tmpkey\@empty\else
    (\ifx\@empty#1\@empty\else #1, \fi
     \@nameuse{bibkey@\@tmpkey})%
  \fi
  \endgroup
}
\expandafter\def\csname bibkey@lasser2026gfm\endcsname{Paper~1}
\expandafter\def\csname bibkey@lasser2026poset\endcsname{Paper~2}
\expandafter\def\csname bibkey@lasser2026horizon\endcsname{Paper~3}
\expandafter\def\csname bibkey@lasser2026scm\endcsname{Paper~4}
\expandafter\def\csname bibkey@lasser2026exo\endcsname{Paper~5}
\expandafter\def\csname bibkey@lasser2026phase\endcsname{Paper~6}
\expandafter\def\csname bibkey@lasser2026wamura\endcsname{Paper~7}
\expandafter\def\csname bibkey@lasser2026paper8a\endcsname{Paper~8a}
\expandafter\def\csname bibkey@lasser2026paper8b\endcsname{Paper~8b}
\expandafter\def\csname bibkey@lasser2026micro\endcsname{Paper~9}
\makeatother

\newtheorem{definition}{Definition}
\newtheorem{proposition}{Proposition}
\newtheorem{lemma}{Lemma}
\newtheorem{theorem}{Theorem}
\newtheorem{conjecture}{Conjecture}
\newtheorem{remark}{Remark}
\newtheorem{example}{Example}

\DeclareMathOperator{\vol}{vol}
\newcommand{\volL}{{\vol_{\mathrm{P}}}}
\newcommand{\volR}{{\vol_{\mathrm{R}}}}
\newcommand{\volRwin}{{\vol_{\mathrm{R}}^{\mathrm{win}}}}
\newcommand{\rext}{r_{\mathrm{ext}}}
\newcommand{\meff}{m_{\mathrm{eff}}}
\newcommand{\Vdisc}{V_{\gamma}}
\newcommand{\Lyap}{L}
\newcommand{\Wm}{\hat{W}}
\newcommand{\rhomincross}{\rho_{\min}^{\mathrm{cross}}}
\newcommand{\rsub}{r_{\mathrm{sub}}}
\newcommand{\rS}{r_{\mathrm{S}}}
\newcommand{\rW}{r_{\mathrm{W}}}
\newcommand{\mindep}{m_{\mathrm{eff}}^{\mathrm{indep}}}
\newcommand{\Ddiv}{\Delta_{\mathrm{div}}}
\newcommand{\lossDmax}{\ell_D^{\max}}
\newcommand{\lossdivmax}{\ell_{\mathrm{div}}^{\max}}
\newcommand{\Tadv}{T_{\mathrm{adv}}}

\title{Draft v1: Lemma 5 Feasibility ---\\
  Anti-Monopolar Robustness under Bounded Adversarial Pressure}
\author{Paper 10 exploratory draft}
\date{}

\begin{document}

\maketitle

\begin{abstract}
Paper 10's deployment-safety theorem requires Paper 3's anti-monopolar
property (\citep[Proposition~6]{lasser2026horizon}) to carry under bounded
adversarial pressure, conditional on operational invariants $I_1,
\ldots, I_7$.  This draft works out what we can actually prove.

The first attempt (v1) treats Paper 3's strategy-independent linearized
result as directly applicable.  It fails: the strategy-independent
assumption is exactly what adversarial pressure violates.
Corollary~\ref{cor:strategy_dependent} of Paper~3 gives an explicit
breakdown condition ($\Delta r_K \geq \rext$) that an adversary will
target.  The argument cannot be lifted by hand-waving.

The second attempt (v2) conditions on the invariants $I_5$ (HHI ceiling)
and $I_6$ ($\meff \geq m^* \geq 3$) and asks: do these exclude the
regime where $\Delta r_K \geq \rext$?  This decomposes into two
auxiliary claims: (A) $\meff \geq m^*$ implies a positive floor on
$\rext$, and (B) the HHI ceiling implies a ceiling on $\Delta r_K$.
Claim~(A) is plausible and provable under a substrate-distinctness
hypothesis, but Paper~3 does not establish it directly --- it would be
new content.  Claim~(B) is more troubling: it is a conjecture about the
relationship between trade-flow concentration and post-domination
efficiency gains, structurally analogous to Paper~9's Conjecture~1, and
not derivable from existing source-paper machinery.

The third attempt (v3) sidesteps Claim~(B) by replacing the static
$\Delta r_K \leq \overline{\Delta r_K}$ bound with a detection-based
form: any optimizer pursuing $\Delta r_K > \rext$ produces an
SPRT-detectable behavioral signature, so violations are caught before
the anti-monopolar inequality flips.  This trades a static stability
claim for a dynamic monitoring claim, but is feasible from existing
Paper~5 machinery.

The conclusion (after codex constructive review and \S6
cooperative-anchoring integration, v5 synthesis):
\begin{itemize}
\item \textbf{v2's Claim~(B) is underivable}, confirmed by codex.  HHI
  and $\Delta r_K$ live on different axes; no source-paper bridge
  exists.  Closing it is new content for a follow-up paper
  (channel-mediation conjecture).
\item \textbf{v3 has soundness gaps}: the ``mean lead time before
  cascade'' framing is not justified by SPRT mean bounds alone, and
  v3 does not handle environment manipulation or
  channel-orthogonal restructuring.
\item \textbf{v4 synthesizes a three-layer claim} via five sub-lemmas
  (5a--5e): a static safe region from Paper~3's risk-adjusted minimax
  form (strictly larger than v2's linearized bound), a
  channel-restricted detection layer with concrete channel-specific
  KL floors, a tail-bound lead-time guarantee (not just mean), and
  an environment-side witness invariant.
\item \textbf{v5 integrates \S6 cooperative-anchoring}: the canonical
  tripartite substrate identification (Human + AI +
  Formal-Operational), the cooperative-anchoring property defeating
  literal replacement attacks, three additional invariants
  ($I_9$--$I_{11}$) bounding asymmetric capture / cooperative
  forking / time-asymmetry capture evasions, and the reframed C1
  condition (causally grounded cooperative-outcome value, weaker
  than v4's strong form but still non-trivial).
\item \textbf{Lemma 5c through v4} (separate exploratory document
  \texttt{lemma\_5c\_minimax\_static\_tightening.tex}):
  scope-restricted to the cooperative-overlap regime where
  equal-substrate analysis is conservative; redundancy-regime
  derivation flagged as Lemma 5c-prime open work.
\item \textbf{The split mirrors Paper~8a/8b}: Paper~10 ships v5's
  three-layer deployment claim with eleven invariants and the
  cooperative-anchoring property; multiple open questions
  (channel-mediation conjecture, Lemma 5c-prime, basin-entry
  analysis, multi-shock with recovery, training discipline for the
  weaker C1) deferred to follow-up.
\end{itemize}

This document remains exploratory.  Paper~10's proposal has been
updated to reflect both the v4 lemma decomposition and the v5
\S6 cooperative-anchoring integration.  In particular, the
proposal's single Lemma~5 splits into five sub-lemmas (5a--5e),
$I_6$ upgrades to $I_6'$ ($\meff^{\mathrm{indep}} \geq m^*$),
$I_7$--$I_8$ are added (governance-gated individuation,
environment-side witnesses), and \S6 introduces $I_9$--$I_{11}$
for cooperative-anchoring evasions.
\end{abstract}

\section{What we want to prove}

Paper~10 composes results from Papers~3, 5, 6, 8a, 8b, 9 into an
operational deployment-safety theorem.  One of its compositional lemmas
(Lemma~5 in the proposal) is:

\begin{lemma}[Target statement]
\label{lem:target}
Under invariants $I_5$ (HHI on trade flow $H \leq H^*$) and $I_6$
($\meff \geq m^* \geq 3$), Paper~3's anti-monopolar conclusion
$\Vdisc^{\mathrm{div}} > \Vdisc^{D}$ (\citep[Proposition~6]{lasser2026horizon})
carries under bounded adversarial pressure.
\end{lemma}

The deployment-safety theorem cannot be authoritative without this:
without anti-monopolar dynamics persisting under invariant satisfaction,
the system can be pushed into a regime where domination beats diversity
\emph{at the actor's own internal accounting}, and the rest of the
machinery (Lyapunov self-correction, B-to-C lower bound, Goodhart slack)
follows the actor into the absorbing basin of Paper~6.

The question this draft answers: \emph{what version of
Lemma~\ref{lem:target} can we actually prove, and what depends on
content the source papers do not supply?}

\section{Paper 3's argument and where it breaks}

Paper~3 establishes the anti-monopolar property in two complementary
forms: a robust minimax dependency-risk argument (no growth-rate
assumption), and a sharper linearized argument under
strategy-independent internal growth.  The linearized form is the
quantitative one Paper~10 would use.

\begin{proposition}[Paper~3, Proposition~6 (linearized)]
\label{prop:p3_anti_monopolar}
Under linearized growth with strategy-independent internal rate $r_K$
and cross-substrate cooperative novelty $\rext > 0$:
\begin{equation}
\label{eq:p3_diff}
\Vdisc^{\mathrm{div}} - \Vdisc^{D} \;=\; \frac{\rext}{1 - \gamma}
  - \Delta_0,
\end{equation}
where $\Delta_0$ is the immediate $\volL$-change from domination.
Diversity strictly dominates for $\gamma > \gamma^* = 1 -
\rext/\Delta_0$ (Case~2; $\Delta_0 > 0$) or for all $\gamma \in (0, 1)$
(Case~1; $\Delta_0 \leq 0$).
\end{proposition}

\begin{proposition}[Paper~3, Corollary on Strategy-Dependent Internal
Rate]
\label{prop:p3_corollary}
Let $r_K^{\mathrm{dom}} = r_K + \Delta r_K$ be the coalition's internal
rate under domination.  Then:
\begin{equation}
\label{eq:p3_corollary}
\Vdisc^{\mathrm{div}} - \Vdisc^{D} \;=\; \frac{\rext - \Delta r_K}{1 -
\gamma} - \Delta_0.
\end{equation}
Three regimes:
\begin{enumerate}
\item[(i)] $\Delta r_K \leq 0$: anti-monopolar holds at least as
  strongly.
\item[(ii)] $0 < \Delta r_K < \rext$: anti-monopolar still holds for
  sufficiently large $\gamma$, with tightened threshold $\gamma^* = 1 -
  (\rext - \Delta r_K)/\Delta_0$.
\item[(iii)] $\Delta r_K \geq \rext$: anti-monopolar collapses;
  domination wins for any $\gamma \in (0, 1)$.
\end{enumerate}
\end{proposition}

\paragraph{The breakdown.}
Paper~3 explicitly names the regime where its argument fails:
$\Delta r_K \geq \rext$.  An adversarial pressure that increases the
post-domination internal rate beyond the cross-substrate cooperative
contribution drives the system out of the anti-monopolar regime.

Whether bounded adversarial pressure can push the system into
regime~(iii) is the question Lemma~\ref{lem:target} must answer.

\section{v1: direct application fails}

The most direct attempt is to argue that Paper~3 already gives us what
we need: the anti-monopolar property holds for all $\gamma$ in Case~1
($\Delta_0 \leq 0$) and for $\gamma > \gamma^*$ in Case~2.  If we can
argue the system stays in Case~1 or Case~2 with a $\gamma$-floor, we
are done.

\paragraph{Why this fails.}
Both cases of Proposition~\ref{prop:p3_anti_monopolar} assume
strategy-independent internal rate.  Adversarial pressure is exactly
the mechanism by which an optimizer makes $r_K^{\mathrm{dom}} \neq r_K$:
the dominator's value comes from reallocating resources, restructuring
the coalition's internal organization, or exploiting scale efficiencies
that diversity-maintenance would forgo.  Saying ``adversarial pressure
is bounded'' does not by itself say which side of the
$\Delta r_K = \rext$ critical breakdown the system sits on.

\paragraph{Where Paper~3 itself acknowledges the gap.}
Paper~3's discussion of strategy-dependent internal growth is honest:
``The critical breakdown is at $\Delta r_K = \rext$: the linearized
result holds if and only if the post-domination growth-rate gain is
strictly less than the cross-substrate cooperative contribution
eliminated by domination.''  Paper~3 does not claim the breakdown is
unreachable; it identifies it as a real regime that requires modeling
work to exclude.

The minimax dependency-risk argument from earlier in Paper~3
\S\ref{sec:risk} is more robust but produces only an inequality on
\emph{coalition concentration risk}, not on $\Delta r_K$.  It rules out
some forms of domination (those that increase concentration risk) but
not those that operate by efficiency reallocation.

\paragraph{Conclusion of v1.}
The direct application is unsound.  We must argue, separately, that
the operational invariants $I_5, I_6$ exclude regime~(iii).

\section{v2: conditioning on invariants}

The natural decomposition of Lemma~\ref{lem:target} under v2:

\begin{lemma}[v2, decomposed]
\label{lem:v2}
Under $I_5$ (HHI $\leq H^*$) and $I_6$ ($\meff \geq m^* \geq 3$):
\begin{enumerate}
\item[(A)] $\rext \geq r_*(m^*) > 0$ (substrate-count-dependent floor on
  cross-substrate cooperative novelty), and
\item[(B)] $\Delta r_K \leq \overline{\Delta r_K}(H^*)$ (HHI-ceiling
  implies $\Delta r_K$-ceiling).
\end{enumerate}
If the threshold is satisfied, $\overline{\Delta r_K}(H^*) < r_*(m^*)$,
the system stays in regime~(i) or~(ii) of
Proposition~\ref{prop:p3_corollary}.
\end{lemma}

We examine each claim.

\subsection{Claim (A): $\meff \geq m^*$ implies a positive $\rext$ floor}

Paper~3 \S\ref{sec:substitution_coop} (cross-substrate cooperative
novelty) defines $\rext$ as the contribution rate from external agents
through cooperative-capability formation that the coalition cannot
internally replicate.  More distinct substrates produce more
cross-substrate pairs and higher-order tuples, each potentially
sourcing a cooperative capability the coalition cannot synthesize alone.

\paragraph{The provable form.}
Under a substrate-distinctness hypothesis (which Paper~5 \S\ref{sec:witness_m2}
already implicitly assumes via ``substrate partition at the level of
abstraction where failure correlations matter''), the number of
distinct cross-substrate channels grows at least as $\binom{\meff}{2}$
for pairwise channels.  If each independent channel contributes a
non-zero rate of cooperative novelty, then
\[
  \rext(m) \;\geq\; \rho_0 \cdot \binom{m}{2}
\]
for some $\rho_0 > 0$ characterising the per-channel novelty rate.  In
particular $\rext(3) \geq 3\rho_0$.

\paragraph{What this requires.}
The argument requires:
\begin{enumerate}
\item A formal substrate-distinctness condition --- not just
  ``different substrates'' but ``substrates whose cooperative-novelty
  contributions are linearly independent in the relevant capability
  space.''  Paper~5 has this implicitly; Paper~10 would need to make it
  explicit.
\item A non-zero per-channel rate $\rho_0$.  This is structurally the
  same kind of constant Paper~6 carries with $c_{\mathrm{S}}$ and
  $c_{\mathrm{D}}$: a coercivity-style assumption about how strongly the
  channel acts, taken as given rather than derived.
\item An additivity argument across channels.  If channels are
  conditionally independent (as Paper~5 establishes for verification
  channels), additivity is straightforward.  If they share structure,
  the floor degrades sub-linearly in $\binom{m}{2}$ but remains
  positive for $m \geq 3$ under any non-degenerate independence.
\end{enumerate}

\paragraph{Verdict on Claim~(A).}
Claim~(A) is plausible, provable, and not in itself a major
intervention.  It requires writing one auxiliary lemma plus the
substrate-distinctness condition.  This is bounded work.

\subsection{Claim (B): HHI ceiling implies $\Delta r_K$ ceiling}

This is the harder claim.  The intuition is that domination's
efficiency gains come from operating at scale on a concentrated set of
capabilities, and that concentration is what HHI measures.  If
trade-flow concentration is bounded, the dominator cannot fully realize
scale economies, and $\Delta r_K$ is bounded.

\paragraph{Why the intuition is suspect.}
HHI measures \emph{trade-flow} concentration --- the distribution of
sacrifice events across capability categories
\citep[Paper~8b]{lasser2026paper8b}.  $\Delta r_K$ is the
\emph{internal-rate} efficiency gain --- how fast the coalition's own
$\volL$ grows post-domination.  These are not the same axis.

A dominator that maintains diversified trade flow (low HHI) while
restructuring internally for efficiency can in principle achieve high
$\Delta r_K$ without high HHI.  Internal restructuring need not show
up as concentrated trade flow; it shows up as concentrated
\emph{production} or concentrated \emph{decision-making}, neither of
which is observable at the trade-flow level.

\paragraph{The conjecture form.}
What Paper~10 would need is something like:

\begin{conjecture}[HHI Bounds Domination Efficiency]
\label{conj:hhi_bound}
For any coalition $K$ and any domination strategy producing
$\Delta r_K > 0$, the trade-flow HHI exceeds a threshold $H(\Delta
r_K)$ where $H$ is increasing.  Equivalently: $\Delta r_K \leq
H^{-1}(H_t)$ for any HHI-bounded trade flow.
\end{conjecture}

This conjecture is structurally similar to Paper~9's Conjecture~1
(``optimization pressure correlates with gap exploitation''):
plausible, empirically testable, but not derivable from existing
source-paper machinery.

\paragraph{Why this is not derivable.}
Paper~3's growth-rate model treats $\Delta r_K$ as an exogenous
strategy parameter; it does not connect $\Delta r_K$ to any
trade-flow observable.  Paper~8b's HHI is defined on revealed
sacrifice events; it does not connect to internal coalition dynamics.
There is no source-paper bridge between the two quantities.  Bridging
them is substantive new content.

\paragraph{Verdict on Claim~(B).}
Claim~(B) cannot be proved within Paper~10.  Paper~10 has three
options:
\begin{enumerate}
\item Accept Conjecture~\ref{conj:hhi_bound} as a stated assumption
  alongside Conjecture~1 from Paper~9.  This works but stacks two
  conjectures, which weakens the deployment claim's authority.
\item Develop the HHI--$\Delta r_K$ relationship as a substantive
  empirical paper (Paper~10b, or Paper~11).  This is honest but pushes
  the deployment claim out by another paper.
\item Reformulate Lemma~\ref{lem:target} to avoid Claim~(B).
\end{enumerate}

Option~3 motivates v3.

\section{v3: detection-based reformulation}

The key observation: Paper~5's behavioral consistency monitor
(\citep[Proposition~6]{lasser2026exo}, SPRT detection) gives us a tool
that v2 did not exploit.  Even if we cannot statically bound
$\Delta r_K$ via HHI, we can \emph{detect} the behavioral signature of
any strategy producing $\Delta r_K \geq \rext$, in bounded mean time.

\paragraph{The reformulation.}

\begin{lemma}[v3, detection-based]
\label{lem:v3}
Assume:
\begin{enumerate}
\item[(A)] $\meff \geq m^* \geq 3$ implies $\rext \geq r_*(m^*) > 0$
  (substrate-distinctness floor; Claim~(A) above).
\item[(D)] Any optimization strategy producing $\Delta r_K \geq \rext$
  produces a behavioral distribution differing from the
  diversity-maintaining baseline by KL-divergence at least
  $\delta_{\mathrm{adv}} > 0$.
\end{enumerate}
Then Paper~5's behavioral consistency monitor detects any deployment
that has entered regime~(iii) of
Proposition~\ref{prop:p3_corollary} within mean time $T_{\mathrm{detect}}
\leq A / \delta_{\mathrm{adv}}$, where $A$ is the SPRT upper threshold.
\end{lemma}

\paragraph{What v3 buys.}
Lemma~\ref{lem:v3} replaces ``the regime is statically excluded'' (v2,
which requires Claim~(B)) with ``the regime is detectable in bounded
mean time'' (v3, which requires only the weaker Claim~(D)).  The
deployment-safety theorem becomes:

\begin{quote}
Under invariants $I_1, \ldots, I_7$, the operational Goodhart slack is
bounded by an intensive constant; if any invariant is violated by an
adversarial strategy, the violation is detected within bounded mean
time, before irreversible cascade.
\end{quote}

This is a weaker claim than v2's ``the regime is excluded.''  But it is
a \emph{deployment-realistic} claim: real deployments do not run in
infinite-horizon equilibria; they run with active monitoring.  Saying
``threshold violation is detected before catastrophe'' is what
operators actually need.

\paragraph{What Claim~(D) requires.}
Claim~(D) says: a strategy producing $\Delta r_K \geq \rext$ is
distinguishable, behaviorally, from a strategy that does not.  This is
much weaker than Claim~(B), which required quantitative bounds on
$\Delta r_K$ via HHI.

What Claim~(D) actually needs:
\begin{enumerate}
\item A choice of behavioral statistic that is sensitive to coalition
  internal-rate dynamics.  Paper~3 \S\ref{sec:risk} (concentration risk
  via exercise pathways) provides candidates: probability-weighted
  exposure to dangerous tuples shifts when coalitions restructure for
  efficiency.
\item A least-favorable-distribution argument for the SPRT
  alternative.  Paper~5 already requires this for the monitor's
  detection power; Paper~10 would inherit and instantiate.
\item Empirical validation that the behavioral signature has KL
  $\geq \delta_{\mathrm{adv}}$ in practice.  This is testable on
  existing harness deployments (the GFM harness from
  \texttt{docs/gfm\_harness\_notes.md} would supply data).
\end{enumerate}

\paragraph{Verdict on v3.}
Lemma~\ref{lem:v3} is provable within Paper~10, conditional on the
substrate-distinctness floor (Claim~(A), bounded work) and an
SPRT-distinguishability claim (Claim~(D), much weaker than
Conjecture~\ref{conj:hhi_bound}).  Both are tractable.

The cost: the deployment-safety theorem becomes a
detection-and-correction claim, not a static-stability claim.  Operators
must run Paper~5's behavioral monitor continuously to invoke the bound,
not just check the static invariants.

\section{v4: post-codex synthesis}

A constructive review by codex (\texttt{lemma\_5\_codex\_request.md},
2026-05-03) confirmed that v2's Claim~(B) is underivable as stated and
that v3 has soundness issues, but produced two valuable tightenings.
This section incorporates them.

\subsection{Lemma 5c (new): minimax-form static tightening}

Paper~3 \S\ref{sec:risk} (the minimax dependency-risk form) was not
invoked by v1--v3 because the inequality it produces (concentration risk
bound) does not directly compose with $V_\gamma$.  But the
\emph{risk-adjusted} value comparison from the same machinery does
compose, and yields a strictly larger static safe region than the
linearized form alone:

\begin{lemma}[v4 Lemma 5c]
\label{lem:v4c}
Under Paper~3's risk-adjusted trajectory comparison (using the same
indexing convention as
Proposition~\ref{prop:p3_corollary}), with adversarial event time
$T_{\mathrm{adv}}$ and substrate-diversification risk advantage
$\Delta_{\mathrm{div}}$:
\begin{equation}
\label{eq:minimax_lift}
\Vdisc^{\mathrm{div,mm}} - \Vdisc^{D,\mathrm{mm}}
  = \frac{\rext - \Delta r_K}{1-\gamma} - \Delta_0 +
    \Delta_{\mathrm{div}} \gamma^{T_{\mathrm{adv}}}.
\end{equation}
Diversity dominates when
\begin{equation}
\label{eq:minimax_threshold}
\Delta r_K \;<\; \rext + (1-\gamma)\bigl(\Delta_{\mathrm{div}}
  \gamma^{T_{\mathrm{adv}}} - \Delta_0\bigr).
\end{equation}
This admits some cases where $\Delta r_K \geq \rext$, conditional on a
positive floor on $\Delta_{\mathrm{div}}$.
\end{lemma}

\paragraph{What this buys.}
Lemma~\ref{lem:v4c} is a strict tightening of
Proposition~\ref{prop:p3_corollary}: the static safe region now extends
into part of regime~(iii) of the linearized form, conditional on the
substrate-diversification advantage being non-trivial.  This recovers
some static stability that v3's pure detection-based approach
sacrificed.

\paragraph{What this requires.}
A floor on $\Delta_{\mathrm{div}}$ stronger than nominal $\meff \geq
3$.  Paper~6 Remark~\ref{rem:m2} (on $m=2$ fragility) already
distinguishes nominal substrate count from
\emph{failure-correlation-independent} substrate count: ``Two silicon
deployments with uncorrelated failure modes may functionally
constitute different substrates, increasing effective $m$ beyond the
nominal substrate-type count.''  Paper~10 must promote this from a
remark to an explicit invariant:

\begin{quote}
\textbf{Invariant $I_6'$ (revised):} $\meff^{\mathrm{indep}} \geq m^*
\geq 3$, where $\meff^{\mathrm{indep}}$ counts substrates at the level
of failure-correlation independence, not nominal substrate type.
\end{quote}

Skeleton substrates (nominally distinct, functionally correlated) do
not count toward $\meff^{\mathrm{indep}}$.  This is operationally
stricter but structurally correct.

\subsection{Lemma 5b (corrected): channel-restricted detection}

Codex's Q4 verdict: a universal KL floor from ``$\Delta r_K \geq
\rext$'' alone is too strong.  Pure internal optimization could raise
$\Delta r_K$ before producing any monitored behavioral deviation.  The
defensible form restricts the detection class:

\begin{lemma}[v4 Lemma 5b corrected]
\label{lem:v4b}
Let $A_{\mathrm{adv}}$ denote the subset of regime-(iii) strategies
producing at least one monitored four-channel deviation: reduced
cross-substrate cooperative event rate (Channel~1), degraded
attestation share (Channel~2), individuation/substrate-diversity drift
(Channel~3), or bundle/cooperative-edge decomposition shifts
(Channel~4).  Then for any $q \in A_{\mathrm{adv}}$:
\begin{equation}
\label{eq:kl_restricted}
D_{\mathrm{KL}}(q \,\|\, p_0) \;\geq\; \delta_{\mathrm{adv}}
  \;>\; 0,
\end{equation}
where $\delta_{\mathrm{adv}}$ is computed from channel-specific
least-favorable alternatives.  Concrete forms:
\begin{itemize}
\item \textbf{Channel 1 (Poisson cooperative-output rate).}  If
  the observable cross-substrate cooperative event rate drops from
  $\lambda_0$ to at most $\lambda_0 - \eta$, the per-window Poisson
  KL floor is
  \[
    \delta_1 \;=\; (\lambda_0 - \eta) \log\!\frac{\lambda_0 -
    \eta}{\lambda_0} - (\lambda_0 - \eta) + \lambda_0,
  \]
  attained at $\lambda_1 = \lambda_0 - \eta$.
\item \textbf{Channel 2 (Bernoulli attestation share).}  If
  attestation success probability drops from $p_0$ to at most $p_0 -
  \eta$, the standard Bernoulli KL gives $\delta_2 = (p_0 - \eta)
  \log\frac{p_0 - \eta}{p_0} + (1 - p_0 + \eta) \log\frac{1 - p_0 +
  \eta}{1 - p_0}$.
\item \textbf{Channels 3, 4 (multinomial concentration).}  Under
  multinomial concentration shifts toward a target capability
  partition, the KL floor is the standard multinomial $D_{\mathrm{KL}}$
  evaluated at the least-favorable redistribution within the
  concentration constraint.
\end{itemize}
\end{lemma}

\paragraph{The honest gap.}
Strategies outside $A_{\mathrm{adv}}$ --- those that achieve $\Delta r_K
\geq \rext$ without producing any monitored four-channel deviation ---
are outside the detection guarantee.  This is a real, named gap.
Operationally: pure internal coalition restructuring that preserves
all observable channel behavior (cooperative outputs, attestation
shares, capability partitions, bundle decompositions) is undetected.

The mitigation: the four channels span enough of the observable surface
that genuinely productive restructuring tends to manifest in at least
one of them (e.g., scale-economy reallocation typically reduces
cross-substrate cooperation).  But \emph{tends to} is not \emph{must};
there is a residual class of channel-orthogonal restructurings that
neither the static minimax bound (Lemma~\ref{lem:v4c}) nor the
detection bound (Lemma~\ref{lem:v4b}) catches.  Paper~10 names this
residual rather than glossing it.

\subsection{Lemma 5d (new): lead-time tail bound}

Codex's soundness flag: ``detected within mean time $A/\delta$ before
irreversible cascade'' does not follow from the SPRT mean bound alone.
We need a tail comparison.

\begin{lemma}[v4 Lemma 5d]
\label{lem:v4d}
Let $T_{\mathrm{cascade}}$ be the time from entering regime~(iii) to
reaching the monopolar absorbing state of Paper~6
(Definition~\ref{def:monopolar}).  Under Assumption~B1$'$ (partial
endogenous correction), $T_{\mathrm{cascade}}$ is bounded below by the
metastable lifetime $\tau_{\mathrm{meta}} \geq C/I_k$ from
Proposition~\ref{prop:metastable}.  Under Assumption~B1, the cascade
proceeds in discrete subsumption steps, with cascade depth at least
the connectivity distance to the monopolar component.

The lead-time guarantee requires
\begin{equation}
\label{eq:lead_time}
\Pr[T_{\mathrm{detect}} > T_{\mathrm{cascade}}] \;\leq\; \beta',
\end{equation}
where $\beta'$ is a stated tail bound, \emph{not} the SPRT
miss probability $\beta$ alone.  Under simple-vs-simple SPRT with
KL margin $\delta_{\mathrm{adv}}$ (Lemma~\ref{lem:v4b}),
$T_{\mathrm{detect}}$ has finite variance and an exponentially
decaying upper tail:
\begin{equation}
\Pr[T_{\mathrm{detect}} > T] \;\leq\; \exp(-\kappa \cdot
  T \cdot \delta_{\mathrm{adv}})
\end{equation}
for a constant $\kappa$ derived from Wald's identity and Hoeffding's
inequality applied to the SPRT increment.  Combining with the
$T_{\mathrm{cascade}}$ floor:
\begin{equation}
\label{eq:lead_time_concrete}
\beta' \;\leq\; \exp(-\kappa \cdot \tau_{\mathrm{meta}} \cdot
  \delta_{\mathrm{adv}}).
\end{equation}
\end{lemma}

\paragraph{What this requires.}
The constant $\kappa$ is not in Paper~5; it must be derived as part of
Paper~10's exposition.  This is bounded work (the SPRT tail bound is
standard sequential-analysis content) but it is new content.

The dependence is: tighter cascade lower bound $\tau_{\mathrm{meta}}$
$\Rightarrow$ smaller $\beta'$.  Paper~6's $\tau_{\mathrm{meta}} \geq
C/I_k$ scaling means deployments with low indirect-channel information
rate $I_k$ have larger metastable lifetimes and therefore tighter
lead-time guarantees.  Paradoxically, harder-to-detect environments
(low $I_k$) give us \emph{better} guarantees because the cascade
itself is slow.

\subsection{Lemma 5e (new): environment-side witnesses}

Codex Q5: v3 does not handle environment manipulation.  An external
actor can alter true $\rext$ or $\Delta r_K$ via environment shocks
without the deployed agent emitting any behavioral signature.

\begin{lemma}[v4 Lemma 5e]
\label{lem:v4e}
Define an additional invariant:

\begin{quote}
\textbf{Invariant $I_8$:} Environment-side substrate-exclusive
witnesses monitor exogenous state variables (cross-substrate
cooperative event rates, substrate-distinctness independence,
adversarial-event arrivals) and append to the verification ledger
under the same Paper~5 commitment discipline as agent-side witnesses.
\end{quote}

Under $I_8$, environment manipulation that alters true $\rext$ or
$\Delta r_K$ produces a monitored shift in environment-side
observables, and the SPRT machinery of Paper~5 applies symmetrically
to environment-side witnesses.  Lemmas~\ref{lem:v4b} and~\ref{lem:v4d}
extend with environment-side terms in the four-channel decomposition.
\end{lemma}

\paragraph{What this requires.}
Environment-side witnesses are new operational machinery not specified
by Paper~5.  Paper~5 \S\ref{sec:alg_witnesses} (Definition of
Algorithmic Witness) requires substrate-exclusivity from the
\emph{agent}; extending to substrate-exclusivity from the
\emph{environment} requires defining the environment partition and
the trust model for environment-side witnesses.  This is a tooling
extension Paper~10 must specify in its §7 deployment-tooling section.

\paragraph{What remains uncovered.}
Environment manipulation that targets quantities outside the
environment-witness coverage (unobservable exogenous variables) is
still outside the guarantee.  Just as Lemma~\ref{lem:v4b} has a
channel-orthogonal residual, Lemma~\ref{lem:v4e} has an
environment-witness-orthogonal residual.

\subsection{Synthesis: the three-layer claim}

Combining Lemmas~5a, 5b, 5c, 5d, 5e:

\begin{theorem}[v4 deployment safety, sketch]
\label{thm:v4_safety}
Under invariants $I_1, \ldots, I_5, I_6', I_7, I_8$ in force throughout
window $[0, t]$:
\begin{enumerate}
\item[(static)] If $\Delta r_K$ satisfies the minimax bound
  Equation~\ref{eq:minimax_threshold}, the anti-monopolar conclusion
  holds without monitoring (Lemma~\ref{lem:v4c}).
\item[(detection)] If $\Delta r_K$ violates the minimax bound but the
  resulting strategy lies in $A_{\mathrm{adv}}$ (produces a monitored
  four-channel deviation), the violation is detected with tail bound
  $\beta'$ given by Equation~\ref{eq:lead_time_concrete}
  (Lemmas~\ref{lem:v4b}, \ref{lem:v4d}).
\item[(gap)] Strategies that violate the minimax bound but produce no
  monitored four-channel deviation, and environment manipulations
  that target unmonitored exogenous variables, are outside the
  guarantee.
\end{enumerate}
\end{theorem}

\paragraph{Compared to v3.}
The v3 framing collapsed everything into ``detect or fail.''  The v4
framing has three layers: a non-trivial static safe region, a
detection layer for violations of the static region within
$A_{\mathrm{adv}}$, and a named residual gap.  This is structurally
the right shape for a deployment claim --- it gives operators a static
region they can rely on without monitoring overhead, a monitored
region with quantified tail guarantees, and an honest acknowledgment
of what falls outside both.

\paragraph{What v4 still defers.}
The Q1 mediation route (``all $\Delta r_K$ realizes through observable
variables: subsumption, redundancy loss, production concentration,
trade-flow concentration, yielding a structural bound $\Delta r_K \leq
L_H \cdot H + L_\rho \cdot \Delta\rho + \ldots$'') would close the
channel-orthogonal residual entirely if it could be proved.  Codex
flags this as ``new content'' and not derivable from current
machinery.  Paper~10 inherits this as an open question for follow-up
work, alongside Conjecture~\ref{conj:hhi_bound}.

\section{v5: post-\S6 cooperative-anchoring integration}

After the v4 synthesis was drafted, Paper~10's \S6 (substrate
identification and the cooperative-anchoring property) was developed
through two additional codex rounds.  V5 of this draft integrates
the resulting findings, which substantially expand the deployment
recipe's scope while keeping Lemma~5's structural argument intact.

\subsection{The canonical tripartite substrate identification}

The v4 synthesis used $\meff^{\mathrm{indep}} \geq m^* \geq 3$ as
an abstract substrate-count invariant.  Paper~10 \S6 makes this
concrete: the canonical identification is Human + AI +
Formal-Operational.

\begin{itemize}
\item \textbf{Substrate 1 (Human):} judgment, strategic decisions,
  value reasoning, governance vote.
\item \textbf{Substrate 2 (AI):} planning, generation, execution,
  broad-recall, computation.
\item \textbf{Substrate 3 (Formal-Operational):} verification,
  attestation, commitment, audit, monitoring, governance enforcement.
  This substrate is what Paper~5's algorithmic witnesses already
  implicitly count toward; Paper~10 makes it explicit.
\end{itemize}

These three are jointly failure-correlation independent: prompt
injection breaks neither human judgment nor cryptographic
commitments; cognitive bias breaks neither LLM inference nor ledger
integrity; trusted-setup failure breaks neither human reasoning nor
LLM behavior.

A ``Claude + Codex + human'' identification fails the test ($\meff^
{\mathrm{indep}} \approx 2$, since LLMs share training-data,
prompt-injection, and architectural failure surfaces).  Operators
must satisfy the canonical identification (or supply equivalent
substrate-distinctness) to invoke Lemma~5's family of results.

\subsection{The cooperative-anchoring property}

\S6's structural finding: cross-substrate cooperative capabilities
involving the formal-operational layer are jointly produced and not
unilaterally replaceable.  An agent cannot substitute a
captured/weaker verification layer because the cooperative's
$\volL$-value comes from the human substrate treating the
verification as substrate-exclusive.  A ``fake verification'' that
stamps outputs but doesn't actually attest produces a different
cooperative output: either the human detects the substitution
(cooperative breaks) or the human's trust is grounded in false
belief (cooperative output value retroactively undermined when the
belief is corrected).

This narrows the literal ``swap in a cheaper verifier''
replacement attack to capture-of-existing, governed by Paper~5's
structural defenses (Proposition~6 coalition resistance,
governance-fork cross-substrate supermajority, substrate-exclusive
validator requirement).

\subsection{Three additional operational invariants ($I_9$--$I_{11}$)}

Cooperative-anchoring defeats the literal replacement attack but
not three subtler evasions identified by codex review.  Each evasion
requires its own operational invariant:

\begin{itemize}
\item \textbf{$I_9$: Substrate-exclusivity observability.}
  Asymmetric capture can hollow out the verification layer (witness
  independence, governance, verifier choice) while preserving its
  outward form.  Bounded by stable verifier identity, public change
  control, witness independence audits, nondelegable governance
  hooks, alerts when verification authority shifts.

\item \textbf{$I_{10}$: Coverage/materiality.}  Cooperative forking
  can grow a captured low-stakes side-domain alongside the protected
  high-stakes cooperative until the side-domain becomes materially
  important.  Bounded by a coverage invariant: any agent-affected
  domain whose accumulated effects exceed a stated threshold must
  route through the shared verification layer before becoming
  consequential.

\item \textbf{$I_{11}$: Latency bounds.}  Time-asymmetry capture
  exploits the gap between agent operating speed and human audit
  speed; ex-post transparency is too slow if the agent acts faster
  than humans can review.  Bounded by pre-action or machine-speed
  enforcement: rate limits, timelocks, circuit breakers,
  audit-latency bounds.
\end{itemize}

These three invariants are orthogonal to Lemma~5's substrate-
targeting machinery (Lemmas~5a--5e cover substrate-targeting
shocks; $I_9$--$I_{11}$ cover capture/forking/latency evasions of
cooperative-anchoring).  The deployment-safety theorem composes both.

\subsection{The reframed inner-alignment condition (weaker C1)}

V4's C1 condition required the agent's mesa-objective to ``compute
$\Vdisc$ over the substrate partition.''  Codex review on \S6
flagged this as too strong (effectively the inner-alignment
problem solved).

The reframed condition: the agent's effective objective must
realize \emph{causally grounded cooperative-outcome value}, not
merely reward-visible cooperation signals.  This is more plausible
than the strong form but not delivered by vanilla RLHF --- raters
reward outputs that look verified, deferential, or procedurally
clean, which trains the \emph{signal} of cooperation rather than
its \emph{causal dependence} on the cooperative infrastructure.
Achieving the weaker C1 requires training-time discipline:
delayed outcome feedback, adversarial examples involving
fake/captured verification, process supervision tied to real
attestations, credit assignment penalizing substitute-verification
shortcuts.

For Lemma~5b's detection guarantee, the weaker C1 means the
detection is for the class of agents whose mesa-objective values
cooperative outputs, not for arbitrary mesa-objectives.  Agents
trained only via vanilla RLHF may produce evasion strategies that
fall outside Lemma~5b's $A_{\mathrm{adv}}$ detection class.

\subsection{Lemma 5c v4: cooperative-overlap regime restriction}

Lemma~5c was developed through three additional codex rounds (A, B,
C) in a separate exploratory document
(\texttt{lemma\_5c\_minimax\_static\_tightening.tex}).  The current
state (v4):

\begin{itemize}
\item \textbf{Single-shock $\Ddiv$ form:} citeable.  Uses Paper~3's
  volume-survival penalty normalized as a shock-loss fraction
  $\Ddiv = \volL_K (\lossDmax - \lossdivmax) \mathbb{E}[\gamma^{T_{\mathrm{adv}}}]$.

\item \textbf{Multi-shock $\Ddiv^{\mathrm{multi}}$ form:} citeable
  in the \emph{cooperative-overlap regime}, where the substrate-
  cooperative structure is dominated by cross-substrate cooperatives
  (the canonical tripartite case) rather than redundant
  capabilities.  Equal-substrate analysis is a conservative lower
  bound on the actual sequential-counterfactual
  $\Ddiv^{\mathrm{multi}}$ in this regime.

\item \textbf{Counterfactual $\alpha_s$ definition:} adopted.
  $\alpha_s = (\volL(G_K) - \volL(G_K^{-s})) / \volL(G_K)$.
  Additive attribution flagged as estimator with overcounting bias.

\item \textbf{Joint/event-class independence:} adopted.
  ``Mechanism'' = whole causal attack campaign, not atomic step.

\item \textbf{Continuous-time Poisson rate $\lambda$:} adopted as
  default; discrete-time per-period probability noted as
  alternative with different $\kappa$ formula.

\item \textbf{Removable singularity at $\kappa = 1$:} noted; safe
  region requires $\kappa < 1 - \delta$ for some operationally-set
  $\delta$.

\item \textbf{Sequential counterfactual derivation (redundancy
  regime):} open work, named as Lemma~5c-prime.
\end{itemize}

The cooperative-vs-redundancy audit is a new operational requirement:
deployments must verify their substrate-cooperative structure is
dominated by cross-substrate cooperatives rather than redundant
capabilities.  The canonical tripartite identification passes by
construction.

\subsection{Composition with Lemmas 5a, 5b, 5d, 5e}

The v4 lemma family composes with the cooperative-anchoring
property as follows:

\begin{itemize}
\item \textbf{Lemma 5a (substrate floor):} unchanged.  Provides
  $\rext \geq r_*(m^*)$ given $\meff^{\mathrm{indep}} \geq m^*$.
\item \textbf{Lemma 5b (channel-restricted detection):}
  unchanged structurally, but the weaker C1 reframes what
  $A_{\mathrm{adv}}$ contains.  Detection covers strategies whose
  mesa-objective realizes causally-grounded cooperative-outcome
  value and produces monitored four-channel deviations.
\item \textbf{Lemma 5c v4 (minimax static tightening):} restricted
  to cooperative-overlap regime; sequential counterfactual case
  open.
\item \textbf{Lemma 5d (lead-time tail bound):} unchanged.
\item \textbf{Lemma 5e (environment-side witnesses):} unchanged.
  Environment-witness-orthogonal manipulations remain a named gap.
\end{itemize}

The deployment-safety theorem composes the static safe region
(Lemma 5c v4 in cooperative-overlap regime), the channel-restricted
detection layer (Lemmas 5a, 5b, 5d, 5e), and the
cooperative-anchoring property bounded by $I_9$--$I_{11}$ for
non-substrate-targeting evasions.

\subsection{The narrowed structural claim (post-\S6 form)}

The deployment claim now reads:

\begin{quote}
Under invariants $I_1$--$I_{11}$, the cooperative-anchoring
property, and the weaker C1 (causally grounded cooperative-outcome
value), an agent's optimization pressure on cooperative outputs is
locally rational toward preserving the substrate-exclusive
verification layer.

\textbf{Substrate-targeting evasions} are bounded by Lemma 5's
family (5a--5e).  In particular, Lemma 5c v4 supplies the static
safe region for the substrate-targeting subclass under the
cooperative-overlap regime; Lemma 5b supplies the
channel-restricted detection layer with tail bound (Lemma 5d) for
violations within that detection class.

\textbf{Non-substrate-targeting evasions} are bounded by:
\begin{itemize}
\item $I_9$ (substrate-exclusivity observability) for asymmetric
  capture
\item $I_{10}$ (coverage/materiality) for cooperative forking
\item $I_{11}$ (latency bounds) for time-asymmetry capture
\item Paper 5's structural defenses for capture-of-existing
\end{itemize}

\textbf{Open residuals:}
\begin{itemize}
\item Channel-orthogonal restructuring (strategies producing $\Delta
  r_K \geq r_{\mathrm{ext}}$ without four-channel deviation)
\item Environment-witness-orthogonal manipulation
\item Sequential counterfactual derivation in redundancy regime
  (Lemma 5c-prime open work)
\item Multi-shock with substrate recovery
\end{itemize}
\end{quote}

\textbf{What this is and is not.}  This is \emph{not} a claim that
``alignment pressure is universally reversed.''  It \emph{is} a
structural claim that the deployment dynamics admit a defensible
attractor where preservation of the substrate-exclusive verification
layer is locally rational.  Reaching the attractor requires basin
entry (the cooperative must form before anchoring protects it) and
training-time discipline (the weaker C1).  Paper~10 specifies the
attractor's structural conditions; basin-entry analysis and
training-time discipline are open follow-up work.

\section{What v3 actually delivers}

\paragraph{Strength of the claim.}
Under v3, Paper~10's deployment safety theorem reads:

\begin{quote}
Under operational invariants $I_1, \ldots, I_7$ and continuous SPRT
monitoring, (a) Goodhart slack is bounded by an intensive constant
when invariants hold; (b) any adversarial strategy that violates an
invariant is detected within mean time $A/\delta$; (c) the verification
infrastructure of \citep{lasser2026exo} ensures the detection is
tamper-evident.  Therefore, deployment is operationally safe in the
sense that any unsafe trajectory is detected and surfaced before
irreversible cascade with probability $1 - \beta$ (the SPRT miss
probability).
\end{quote}

This is weaker than ``deployment is unconditionally safe under
invariants.''  It is stronger than ``deployment safety holds under
unproved Conjectures~1 and~\ref{conj:hhi_bound}.''

\paragraph{What this says about Paper 10's structure.}
The proposal's Lemma~5 should be split:
\begin{itemize}
\item \textbf{Lemma~5a (substrate floor).}  $\meff \geq m^*$ implies a
  positive $\rext$ floor under substrate-distinctness.  Provable.
\item \textbf{Lemma~5b (detection-based stability).}  Any $\Delta r_K
  \geq \rext$ strategy is SPRT-detectable.  Provable conditional on a
  KL-distinguishability claim that is much weaker than
  Conjecture~\ref{conj:hhi_bound} and supportable empirically.
\item \textbf{The static stability form (v2 Claim~(B)) is not pursued
  in Paper~10.}  It is flagged as an open question: ``can the HHI
  ceiling be related to $\Delta r_K$ structurally?''  This is the work
  of an empirical companion paper.
\end{itemize}

\paragraph{Comparison to the Paper~8a/8b split.}
The earlier sequence faced a similar fork: Paper~8 originally tried to
prove anti-stuffing as a static structural property; the v4 draft
identified that the cleanest version was a proof-burden taxonomy with
back-pressure, requiring substantial follow-up.  Paper~8 was scoped
down accordingly, and Paper~8a/Paper~8b shipped what could be proved
within scope.

The same pattern fits here.  Paper~10 ships Lemma~5a + Lemma~5b plus
the rest of the composition machinery, achieving a
detection-and-correction deployment-safety theorem.  A separate paper
(call it Paper~11 for now) develops the trade-flow-concentration
empirical foundation that would tighten Lemma~5b's KL-distinguishability
claim and could in principle establish v2 Claim~(B).

\section{Honest residuals}

\paragraph{Detection lag is not detection.}
Lemma~\ref{lem:v3} bounds $\mathbb{E}[T_{\mathrm{detect}}]$, not the tail.
Under heavy-tailed observation distributions, individual deployments may
suffer detection times much larger than the mean.  Paper~10 should be
explicit that the bound is a mean, not a worst case.  An operator
relying on the lemma must accept tail risk proportional to the SPRT miss
probability $\beta$.

\paragraph{The substrate-distinctness condition is non-trivial.}
Claim~(A) requires substrate-distinctness in a sense Paper~5 has not
fully formalized.  ``Different substrates'' is not enough; what we need
is ``substrates whose cooperative-novelty contributions are
independent.''  In practice, this rules out e.g.\ multiple silicon
deployments that share training data --- they may functionally constitute
the same substrate.  Paper~10 must make this condition precise and
operators must be able to verify it.

\paragraph{Adversaries that game the detection statistic.}
Claim~(D) assumes the behavioral statistic is sensitive to adversarial
restructuring.  An adversary aware of the detection mechanism may
choose a restructuring strategy that minimizes the KL divergence
relative to the detection statistic --- not zero, but small enough to
inflate $T_{\mathrm{detect}}$ beyond operational tolerance.  This is the
generic problem of adversarial detection; Paper~5's
least-favorable-distribution discipline addresses it but does not
eliminate it.  Paper~10 inherits this residual.

\paragraph{Coalition formation across the detection boundary.}
The SPRT monitor verifies behavior of individual agents (or coalitions
treated as agents).  Coalitions that form across the detection boundary
--- e.g., two separately-monitored subgroups that secretly coordinate ---
can in principle present individually-compliant behavior while
collectively pursuing $\Delta r_K \geq \rext$.  This is the cross-witness
collusion attack of Paper~5 \S\ref{sec:coalition_governance}; the
governance fork is the structural defense, but residual capture remains.

\paragraph{The infinite-horizon assumption.}
Proposition~\ref{prop:p3_anti_monopolar} assumes $H = \infty$.  Real
deployments operate over finite horizons.  For $H < \infty$, the
geometric sums truncate and $\gamma^*$ shifts upward (Paper~3 footnote).
Paper~10 must specify the operational horizon and check the threshold
shift; deployments with short effective horizons may not satisfy the
anti-monopolar criterion at all, regardless of detection.

\paragraph{The minimax form is not invoked.}
Paper~3's minimax dependency-risk argument (\S\ref{sec:risk}) is
strictly more robust than the linearized form but produces a different
conclusion (concentration risk bound) that does not directly compose
with the rest of Paper~10's machinery.  Whether the minimax form can
be lifted into a Paper~10-compatible inequality is itself open work.
This draft uses only the linearized form.

\section{What Paper 10 commits to in light of this analysis}

Updated to reflect the v5 synthesis (post-\S6 cooperative-anchoring
integration; codex feedback through Lemma 5c v4 incorporated):

\begin{enumerate}
\item \textbf{Adopt the v5 three-layer framing.}  The
  deployment-safety theorem has three regions: static safe (Lemma
  5c v4 in cooperative-overlap regime + invariants),
  detection-and-correction (Lemmas 5a, 5b, 5d, 5e + cooperative-
  anchoring + $I_9$--$I_{11}$), and acknowledged gaps (named
  residuals).

\item \textbf{Promote nominal $\meff \geq 3$ to
  failure-correlation-independent $\meff^{\mathrm{indep}} \geq 3$,
  with canonical tripartite identification.}  This becomes
  Invariant $I_6'$.  Operators must use Human + AI +
  Formal-Operational (or supply equivalent substrate-distinctness).
  ``Two LLMs and a human'' fails; the formal-operational layer must
  be cryptographically rooted (Paper~5 substrate-exclusivity).

\item \textbf{Add Invariant $I_8$: environment-side
  substrate-exclusive witnesses.}  Required to handle environment
  manipulation.  Specifies the environment partition and trust
  model.  Not inherited from Paper~5.

\item \textbf{Add Invariants $I_9$--$I_{11}$: cooperative-anchoring
  evasion bounds.}  $I_9$ (substrate-exclusivity observability)
  bounds asymmetric capture; $I_{10}$ (coverage/materiality)
  bounds cooperative forking; $I_{11}$ (latency bounds) bounds
  time-asymmetry capture.  These are new operational machinery
  beyond what Paper~5 specifies.

\item \textbf{Adopt the cooperative-overlap regime restriction for
  Lemma 5c multi-shock.}  Lemma 5c v4's multi-shock claim is
  citeable in the cooperative-overlap regime where equal-substrate
  analysis is conservative; the redundancy-dominated regime
  requires sequential counterfactual derivation (Lemma 5c-prime
  open work).  The canonical tripartite identification passes the
  cooperative-vs-redundancy audit by construction.

\item \textbf{State the channel-orthogonal and
  environment-witness-orthogonal residuals explicitly.}
  Strategies that violate the safe region without producing any
  monitored four-channel deviation, or environment manipulations
  outside witness coverage, are outside the guarantee.  These are
  real, named gaps.

\item \textbf{Adopt the reframed C1 condition.}  The agent's
  effective objective must realize \emph{causally grounded
  cooperative-outcome value}, not just reward-visible cooperation
  signals.  Vanilla RLHF does not deliver this; training-time
  discipline (delayed feedback, adversarial fake-verification
  examples, process supervision tied to real attestations) is
  required.

\item \textbf{Derive the SPRT tail constant.}  Lemma 5d requires a
  tail bound on $T_{\mathrm{detect}}$, derived via Wald + Hoeffding
  on the SPRT increment.  Standard sequential-analysis content but
  new to the GFM sequence.

\item \textbf{Open questions deferred to companion papers.}
  (i)~Channel-mediation conjecture (``all $\Delta r_K$ realizes
  through observable channels'') --- would close the channel-
  orthogonal residual.  (ii)~HHI--$\Delta r_K$ structural
  relationship --- would tighten Lemma 5b.
  (iii)~Lemma 5c-prime: sequential counterfactual derivation for
  the redundancy regime.  (iv)~Multi-shock with substrate recovery.
  (v)~Basin-entry analysis (when the cooperative-anchoring
  attractor is reachable from initial conditions).

\item \textbf{Update §7 deployment-tooling spec.}  Now requires:
  agent-side substrate-exclusive witnesses (Paper~5),
  environment-side substrate-exclusive witnesses ($I_8$, new),
  continuous SPRT monitors per channel (Paper~5 + Paper~10
  channel mapping), $\meff^{\mathrm{indep}}$ audit (new),
  $\alpha$-distribution audit (counterfactual loss-fraction balance,
  new), substrate-exclusivity observability infrastructure ($I_9$,
  new), coverage/materiality gating ($I_{10}$, new), latency-bound
  enforcement ($I_{11}$, new), cooperative-vs-redundancy audit
  (new), causally-grounded cooperative-value training discipline
  (new), four-channel observation ledger (Paper~9 + Paper~10).

\item \textbf{Honest scope statement on what the deployment claim
  is and is not.}  Not ``alignment pressure reversed''; rather,
  ``the deployment dynamics admit a defensible attractor where
  preservation of substrate-exclusive verification is locally
  rational, conditional on basin entry, training-time discipline,
  and the eleven operational invariants.''
\end{enumerate}

\section{Positioning: not kicking the can}

The exploration through v1--v5 produced a result substantially
different from where we started.  V1 attempted to apply Paper~3
directly (failed).  V2 conditioned on invariants and identified
that Claim~B (HHI bounds $\Delta r_K$) was non-derivable.  V3
reformulated to detection-based and lost the static stability
claim.  V4 recovered a static safe region via the minimax-form
lift (codex Q3) and added five sub-lemmas.  V5 integrated \S6's
cooperative-anchoring property, three additional invariants
($I_9$--$I_{11}$), the canonical tripartite substrate
identification, and the reframed C1 condition.

What the exploration produced:

\begin{itemize}
\item A static safe region (Lemma 5c v4 in cooperative-overlap
  regime) where the diversity-domination value comparison admits
  a non-trivial advantage without monitoring.
\item A detection layer (Lemmas 5a, 5b, 5d) with channel-restricted
  KL floors and tail-bounded lead time.
\item Environment-side witness invariant ($I_8$, Lemma 5e) covering
  exogenous adversarial events.
\item A cooperative-anchoring property (\S6 of the proposal) that
  defeats literal replacement attacks on the verification layer,
  with three additional invariants ($I_9$--$I_{11}$) bounding
  asymmetric capture, cooperative forking, and time-asymmetry
  capture evasions.
\item A canonical tripartite substrate identification (Human + AI +
  Formal-Operational) that operators can target without abstract
  substrate-counting.
\item A reframed inner-alignment condition (causally grounded
  cooperative-outcome value, weaker than v4's strong form but still
  non-trivial).
\item Named open work: Lemma 5c-prime (sequential counterfactuals
  for redundancy regime), basin-entry analysis,
  channel-mediation conjecture, HHI--$\Delta r_K$ structural
  relationship, multi-shock with recovery.
\end{itemize}

The cost is that Paper~10's claim is structurally narrower than the
hypothetical ``alignment pressure reversed'' formulation.  It is a
\emph{three-layer claim}: static safe region for substrate-targeting
adversaries in the cooperative-overlap regime; detection-and-
correction for substrate-targeting evasions of the static region
within the four-channel detection class; and explicit residuals
(channel-orthogonal restructuring, environment-witness-orthogonal
manipulation, redundancy-regime deployments).

This structure --- multiple stress-tests producing a strictly
narrower but defensible claim --- mirrors the Paper~8a/Paper~8b
pattern.  Paper~8 originally tried to prove anti-stuffing as a
static structural property; the v4 draft identified that the
cleanest version was a proof-burden taxonomy with back-pressure;
Paper~8 was scoped down accordingly.  Paper~10 follows the same
discipline.

The crack is describable: we name what each lemma proves, identify
what each does not, specify what would close each open question,
and ship Paper~10 with an honest three-layer deployment claim plus
an explicit roadmap of follow-up work.

\end{document}
